% Sinh tự động — scripts/make_report_data.py
\begin{table}[htbp]
  \centering
  \scriptsize
  \setlength{\tabcolsep}{4.5pt}
  \begin{tabular}{@{}l r r r r r r@{}}
    \toprule
    & \multicolumn{2}{c}{10\,\% nhãn} & \multicolumn{2}{c}{20\,\% nhãn} & \multicolumn{2}{c}{50\,\% nhãn} \\
    \cmidrule(lr){2-3}\cmidrule(lr){4-5}\cmidrule(lr){6-7}
    Cấu hình & Macro-F1 & $\Delta$ & Macro-F1 & $\Delta$ & Macro-F1 & $\Delta$ \\
    \midrule
    \cfg{E1} & 0,9319 & --- & 0,9497 & --- & 0,9690 & --- \\
    \addlinespace[2pt]
    \cfg{E2-L} & 0,9313 & $-$0,06 & 0,9505 & +0,08 & 0,9745 & +0,55 \\
    \addlinespace[2pt]
    \cfg{E2} & 0,9470 & +1,51 & 0,9631 & +1,34 & 0,9777 & +0,88 \\
    \cfg{E3} & 0,9438 & +1,19 & --- & --- & --- & --- \\
    \cfg{E4} & 0,9527 & +2,07 & --- & --- & --- & --- \\
    \cfg{E5} & 0,9496 & +1,77 & --- & --- & --- & --- \\
    \cfg{E6} & 0,9579 & +2,60 & --- & --- & --- & --- \\
    \bottomrule
  \end{tabular}
  \caption{Kết quả trên $\Dtest$, trung bình qua ba lần lặp. $\Delta$ là chênh lệch Macro-F1 ghép cặp so với \cfg{E1} ở cùng mức nhãn, đơn vị điểm phần trăm. Hai nhóm được ngăn bằng khoảng trắng: \cfg{E1}/\cfg{E2-L}/\cfg{E2} là khối chính; \cfg{E3}--\cfg{E6} là khối mở rộng contrastive.}
  \label{tab:extension-results}
\end{table}